\section{Background}
\label{sec:background}

\subsection{The deployment, and what a principal reaches}
\label{sec:model}

A deployment is a tuple $D = (P, A, G, T)$.

\begin{description}
\itemsep1pt
\item[$P$] the principals: every role and user the deployment declares.
\item[$A$] the agent runtimes. Each $a \in A$ carries an artifact and an
  execution role $\rho(a) \in P$.
\item[$G \subseteq P \times \mathit{Act} \times \mathit{Res}$] the grants,
  taken from the policies the deployment attaches. $(p, \alpha, r) \in G$
  means $p$ may perform action $\alpha$ on resource $r$.
\item[$T \subseteq P \times P$] the trust relation. $(p,q) \in T$ when $q$'s
  trust policy admits $p$.
\end{description}

What the deployment authorised $p$ to reach is read straight off the grants:
%
\begin{equation}
\label{eq:auth}
\auth(p) = \{\, (\alpha, r) : (p, \alpha, r) \in G \,\}.
\end{equation}
%
Write $E(p)$ for the set of principals $p$ can cause code to run as. What $p$
reaches in effect is the least fixpoint of
%
\begin{equation}
\label{eq:eff}
\eff(p) \;=\; \auth(p) \;\cup \bigcup_{q \in E(p)} \eff(q),
\end{equation}
%
and an escalation exists when the two differ:
%
\begin{equation}
\label{eq:esc}
\Esc(p) \iff \eff(p) \setminus \auth(p) \neq \emptyset .
\end{equation}
%
$E(p)$ is what a single-identity analysis leaves out. It computes $\auth(p)$
and reports on it, and the difference in Equation~\ref{eq:esc} needs the whole
deployment rather than one principal at a time.
Section~\ref{sec:approach} gives the three rules that produce $E(p)$ on this
platform.

\subsection{Agent runtimes on \agentcore}

An \agentcore agent runtime is an AWS resource carrying the container image
or code archive that holds the agent, environment variables, a network
configuration, and an execution role. Two properties of it matter here.

\paragraph{One call changes everything about a runtime.}
\api{Update\allowbreak AgentRuntime} takes the artifact and the role ARN, and the API
reference marks both as required~\cite{awsAgentCoreUpdateRuntime}. A caller that holds the permission to update
a runtime therefore decides three things at once: the code that will run, the
identity it runs as, and, through the environment variables in the same
request, the data it is pointed at.

A published route reaches the same role without this call, by running a
command inside an existing runtime~\cite{garcia2026agentcorepaths}, leaving
the runtime unchanged.

\paragraph{Credentials arrive out of band.}
The runtime reads its role credentials from instance metadata rather than from
an environment variable, and the only control over that channel is
\api{metadata\allowbreak Configuration.\allowbreak requireMMDSV2}. Code that
replaces the agent's own code inherits the channel without carrying a
credential.

\subsection{\iam{iam:PassRole} and the managed policy}

\iam{iam:PassRole} is the permission to hand a role to a
service~\cite{awsPassRole}. AWS ships a managed policy for this platform,
\iam{BedrockAgentCoreFullAccess}~\cite{awsBedrockAgentCoreFullAccess}. The version analysed here is v21, read on 8 November 2026. Its
\iam{iam:PassRole} statement is scoped by two conditions: the role name must
match \iam{arn:aws:iam::\allowbreak *:role/\allowbreak *BedrockAgentCore*}, and the role must be passed
to \iam{bedrock-agentcore.amazonaws.com}.

The first condition is a name match. It does not distinguish one agent's
execution role from another's, and the roles the AWS console creates by default
for this service satisfy the pattern. Section~\ref{sec:case-study} shows what
follows from that.

\subsection{Trust policies in the vendor template}

The trust policy the documentation gives for an agent execution role scopes
\iam{sts:AssumeRole} with \iam{aws:SourceArn} set to
\iam{arn:aws:\allowbreak bedrock-agentcore:\allowbreak REGION:ACCOUNT:*}. That
wildcard covers every agent runtime in the account.

\subsection{Connections on \foundry}

\foundry has a different shape. An account, and each project inside it, holds
\emph{connections}, and a connection records how the account authenticates to
the resource it reaches. The \api{authType} field~\cite{msFoundryConnections} decides whether the call
presents a directory identity, such as \api{AAD} or \api{ManagedIdentity}, or a
credential stored in the connection, such as \api{ApiKey} or \api{SAS}.

Where several agents share a project, a connection holding a stored credential
is presented identically by all of them. The target cannot tell the agents
apart in its logs, and each agent reaches everything the connection is trusted
with. Section~\ref{sec:foundry} measures how common each choice is.
